\section{Cele i scenariusze testowe}

Celem badań eksperymentalnych jest porównanie dwóch wariantów architektonicznych tej samej aplikacji: wariantu monolitycznego oraz wariantu mikroserwisowego. Badanie zostało zaprojektowane w taki sposób, aby ocenić wpływ stylu architektonicznego na trzy główne obszary analizy: wydajność, skalowalność oraz koszty utrzymania w środowisku chmurowym.
\par
W ramach eksperymentu przyjęto kilka scenariuszy testowych, różniących się konfiguracją środowiska oraz sposobem skalowania systemu. Pozwala to ocenić nie tylko zachowanie aplikacji przy standardowym obciążeniu, lecz także jej reakcję na wzrost zapotrzebowania na zasoby. Każdy scenariusz został przygotowany tak, aby umożliwić bezpośrednie porównanie wyników uzyskanych dla obu badanych architektur.
\par
Scenariusze testowe obejmują trzy podstawowe grupy eksperymentów. Pierwsza z nich dotyczy testów bez zastosowania mechanizmów skalowania i pozwala ocenić bazową wydajność obu wariantów. Druga grupa obejmuje testy skalowania wertykalnego, w których zwiększane są zasoby pojedynczych instancji aplikacji. Trzecia grupa obejmuje testy skalowania horyzontalnego, polegające na zwiększaniu liczby instancji systemu lub jego wybranych komponentów.
\par
Tak zdefiniowany układ eksperymentów umożliwia ocenę zachowania aplikacji w różnych warunkach pracy i stanowi podstawę do dalszej analizy porównawczej.

\section{Narzędzia wykorzystane do pomiarów}

Przeprowadzenie badań eksperymentalnych wymagało użycia zestawu narzędzi umożliwiających generowanie obciążenia, monitorowanie działania aplikacji, obserwację wykorzystania zasobów oraz gromadzenie wyników pomiarów. Zastosowane narzędzia dobrano w taki sposób, aby zapewnić powtarzalność testów oraz możliwość porównania wyników pomiędzy kolejnymi uruchomieniami eksperymentu.
\par
Do generowania ruchu testowego wykorzystano narzędzie służące do przeprowadzania testów obciążeniowych, które umożliwia definiowanie scenariuszy żądań, kontrolę poziomu równoległości oraz pomiar podstawowych parametrów wydajnościowych, takich jak czas odpowiedzi i liczba obsłużonych żądań. Narzędzie to stanowi podstawę do oceny zachowania systemu pod rosnącym obciążeniem.
\par
Do monitorowania środowiska wykorzystano mechanizmy pozwalające na obserwację parametrów infrastrukturalnych, takich jak zużycie procesora, pamięci operacyjnej, ruch sieciowy oraz liczba aktywnych instancji aplikacji. Dane te są niezbędne do oceny efektywności wykorzystania zasobów oraz do analizy wpływu skalowania na zachowanie systemu.
\par
W badaniu wykorzystano również narzędzia wspierające automatyzację wdrożeń i wykonywania eksperymentów, co pozwoliło ograniczyć udział działań manualnych i zwiększyć porównywalność wyników. Dzięki temu każda seria testowa mogła zostać przeprowadzona według tej samej procedury.

\section{Metryki oceny}

Ocena badanych architektur została oparta na zestawie metryk pozwalających uchwycić zarówno wydajność aplikacji, jak i konsekwencje infrastrukturalne oraz kosztowe zastosowania danego stylu architektonicznego. Dobór metryk wynika bezpośrednio z celu pracy i podporządkowany jest trzem głównym obszarom porównania: wydajności, skalowalności oraz kosztom utrzymania.

\subsection{Przepustowość}

Przepustowość określa liczbę żądań, które system jest w stanie obsłużyć w jednostce czasu. Jest to jedna z podstawowych metryk oceny wydajności aplikacji, szczególnie istotna w testach obciążeniowych. W kontekście niniejszego badania przepustowość pozwala określić, który wariant architektoniczny efektywniej przetwarza rosnący strumień żądań.
\par
Analiza przepustowości jest szczególnie ważna przy porównywaniu systemów wdrożonych w środowisku chmurowym, ponieważ pozwala ocenić, czy wzrost wykorzystania zasobów przekłada się na rzeczywisty wzrost zdolności przetwarzania. W badaniu przepustowość stanowi jedną z głównych metryk używanych do porównania architektury monolitycznej i mikroserwisowej.

\subsection{Latencja}

Latencja opisuje opóźnienie występujące pomiędzy wysłaniem żądania a otrzymaniem odpowiedzi przez klienta. W praktyce jest ona zwykle wyrażana jako czas odpowiedzi systemu i stanowi kluczową metrykę z punktu widzenia jakości odbieranej przez użytkownika końcowego.
\par
W niniejszej pracy analiza latencji pozwala ocenić wpływ rozproszenia logiki biznesowej i komunikacji pomiędzy usługami na czas realizacji żądań. Ma to szczególne znaczenie przy porównywaniu systemu monolitycznego z wariantem mikroserwisowym, w którym dodatkowy narzut może wynikać z wywołań sieciowych, serializacji danych oraz pośrednictwa pomiędzy komponentami.
\par
W badaniu analizowane są zarówno wartości średnie, jak i wybrane statystyki pozycyjne czasu odpowiedzi, co pozwala pełniej ocenić zachowanie aplikacji pod obciążeniem.

\subsection{Zużycie zasobów obliczeniowych}

Zużycie zasobów obliczeniowych obejmuje przede wszystkim wykorzystanie procesora, pamięci operacyjnej oraz, w razie potrzeby, ruchu sieciowego i innych parametrów infrastrukturalnych. Metryka ta ma istotne znaczenie przy ocenie efektywności działania systemu, ponieważ pozwala określić, jakim kosztem infrastrukturalnym osiągane są określone wyniki wydajnościowe.
\par
W kontekście porównania dwóch stylów architektonicznych analiza zasobów umożliwia ocenę, czy większa modularność i elastyczność mikroserwisów nie prowadzi jednocześnie do nadmiernego narzutu operacyjnego. W przypadku architektury monolitycznej metryka ta pozwala z kolei ocenić, czy prostszy model uruchomieniowy rzeczywiście przekłada się na bardziej ekonomiczne wykorzystanie infrastruktury.

\subsection{Koszty utrzymania}

Koszty utrzymania stanowią czwarty wymiar oceny badanych architektur i obejmują zarówno koszty bezpośrednie związane z wykorzystaniem zasobów chmurowych, jak i koszty pośrednie wynikające ze złożoności środowiska. W praktyce analiza ta uwzględnia przede wszystkim koszt instancji obliczeniowych, zasobów pamięciowych, mechanizmów równoważenia obciążenia, transferu danych oraz komponentów wspierających działanie systemu.
\par
W przypadku niniejszego badania analiza kosztów nie ogranicza się wyłącznie do nominalnej ceny zasobów, lecz obejmuje także relację pomiędzy ponoszonym kosztem a uzyskiwaną wydajnością i skalowalnością. Taki sposób oceny pozwala lepiej określić rzeczywistą opłacalność zastosowania danego stylu architektonicznego.

\section{Scenariusz testów bez skalowania}

Pierwszy scenariusz eksperymentalny obejmuje testy przeprowadzane bez zastosowania dodatkowych mechanizmów skalowania. Celem tego etapu jest określenie bazowych właściwości obu wariantów aplikacji przy stałej konfiguracji środowiska uruchomieniowego.
\par
W tym scenariuszu zarówno wariant monolityczny, jak i mikroserwisowy uruchamiane są w porównywalnej konfiguracji zasobów. Następnie do obu systemów kierowany jest ruch testowy o stopniowo rosnącym natężeniu, co pozwala zaobserwować momenty pogorszenia wydajności, wzrostu opóźnień oraz zmian w wykorzystaniu zasobów.
\par
Scenariusz ten pełni rolę punktu odniesienia dla dalszych eksperymentów związanych ze skalowaniem i umożliwia ocenę, jak obie architektury zachowują się w warunkach bazowych.

\section{Scenariusz testów skalowania wertykalnego}

Drugi scenariusz obejmuje badanie wpływu skalowania wertykalnego na zachowanie systemu. Skalowanie wertykalne polega na zwiększaniu zasobów pojedynczej instancji aplikacji lub usługi, na przykład poprzez przydzielenie większej liczby rdzeni procesora lub większej ilości pamięci operacyjnej.
\par
W ramach tego scenariusza wykonywane są kolejne serie testów dla różnych konfiguracji zasobów, przy zachowaniu tego samego schematu obciążenia. Pozwala to ocenić, w jakim stopniu zwiększenie parametrów pojedynczej instancji przekłada się na poprawę przepustowości, obniżenie latencji oraz zmianę kosztów utrzymania.
\par
Scenariusz ten jest szczególnie istotny dla architektury monolitycznej, która często opiera się właśnie na skalowaniu pionowym, ale umożliwia również ocenę, w jakim stopniu taki mechanizm jest efektywny dla wybranych komponentów systemu mikroserwisowego.

\section{Scenariusz testów skalowania horyzontalnego}

Trzeci scenariusz dotyczy skalowania horyzontalnego, polegającego na zwiększaniu liczby instancji aplikacji lub wybranych usług. W środowiskach chmurowych jest to jeden z podstawowych mechanizmów zwiększania zdolności systemu do obsługi rosnącego obciążenia.
\par
W przypadku wariantu monolitycznego skalowanie horyzontalne obejmuje replikację całej aplikacji, natomiast w wariancie mikroserwisowym możliwe jest niezależne zwiększanie liczby instancji wybranych usług. Pozwala to zbadać, czy większa elastyczność architektury mikroserwisowej przekłada się na bardziej efektywne wykorzystanie zasobów i lepsze wyniki przy rosnącym obciążeniu.
\par
W ramach eksperymentu analizowany jest wpływ liczby instancji na przepustowość, czas odpowiedzi, zużycie zasobów oraz koszt działania systemu. Wyniki tego scenariusza pozwalają ocenić praktyczną skuteczność skalowania obu badanych architektur.

\section{Sposób gromadzenia i przetwarzania wyników}

Dane pomiarowe uzyskane w trakcie eksperymentów są gromadzone dla każdej serii testowej w sposób uporządkowany i umożliwiający ich dalszą analizę porównawczą. Rejestrowane są zarówno metryki wydajnościowe pochodzące z narzędzi testowych, jak i dane infrastrukturalne związane z wykorzystaniem zasobów oraz konfiguracją środowiska.
\par
W celu zwiększenia wiarygodności wyników każda seria testowa może być wykonywana wielokrotnie, a do dalszej analizy wykorzystywane są wartości uśrednione lub reprezentatywne statystyki opisowe. Takie podejście pozwala ograniczyć wpływ chwilowych odchyleń wynikających z niestabilności środowiska chmurowego lub zakłóceń sieciowych.
\par
Przetwarzanie wyników obejmuje ich agregację, porządkowanie oraz przygotowanie do prezentacji w formie wykresów, tabel i zestawień porównawczych. Na tej podstawie możliwa jest dalsza interpretacja zachowania badanych architektur w różnych scenariuszach testowych.

\section{Ograniczenia metodyki badawczej}

Mimo przyjęcia uporządkowanej i możliwie powtarzalnej metodyki, przeprowadzone badanie posiada pewne ograniczenia. Po pierwsze, eksperyment opiera się na jednej analizowanej aplikacji, co oznacza, że uzyskane wyniki nie muszą być bezpośrednio uogólnialne na wszystkie klasy systemów informatycznych.
\par
Po drugie, na wyniki eksperymentu mogą wpływać cechy konkretnego środowiska chmurowego, w tym sposób przydzielania zasobów, chwilowe obciążenie infrastruktury współdzielonej oraz właściwości sieci. Czynniki te są częściowo ograniczane przez powtarzalność testów, jednak nie mogą zostać całkowicie wyeliminowane.
\par
Po trzecie, przyjęte scenariusze testowe koncentrują się na wybranym punkcie końcowym i określonym modelu obciążenia. Oznacza to, że inne typy żądań lub inna struktura ruchu mogłyby prowadzić do częściowo odmiennych obserwacji. Mimo tych ograniczeń przyjęta metodyka umożliwia przeprowadzenie spójnego i użytecznego porównania badanych architektur.